\section{Relational coding of the diagonal argument}
\label{rel-section}

We continue with the fixed labelled wrappers, their data cancellation law
(Lemma~\ref{lem:enc-data-cancellation}), and their predicate evaluation laws
\eqref{chan-predicate-evaluation} and \eqref{chan-x-evaluation}.
In particular, for either of the
predicate spaces under consideration,
\[
 \operatorname{pred}(y.\Phi)@w
       \ \leftrightarrow\ \Phi[\operatorname{Round}w/y].
\]
Until a specialization is explicitly made, \(A:s\) and \(m:M(A)\) are
open data parameters.  Write \(L=L(A)\).  The variables \(y,z,w,a\)
range over \(L\), \(P,P'\) over \(\mathsf P_L\), \(x,x'\) over \(X\),
and \(i,i'\) over \(\mathsf P_X\).  The spaces \(V,X,\mathsf P_X\)
are independent of \(A,m\).  All displayed quantifiers use the fixed
direct-vertex annotations: \(j\) for \(L,X\) and \(h\) for
\(\mathsf P_L,\mathsf P_X\).  The other parameters are parameters of
term builders, not additional logical quantifiers.

The operations \(\mathsf{pl}\), \(\mathsf{All}\), \(\mathsf{Ex}\),
\(\mathsf{Gd}\), \(\operatorname{to}\), and
\(\operatorname{back}\) have the typed expansions already constructed.
For clarity, when a proof below says that specified proofs yield a
formula by \(\mathsf{pl}\), that formula is its target.  Quantifiers
are treated as propositional atoms until explicitly introduced or
instantiated.  An existential is eliminated only into the displayed
target, which is independent of its locally opened witnesses.

Our goal is a proof of the terminal formula $F_b$ in the ambient labelled
context. First we construct relations and probes for open parameters
$A,m$. Next we choose $A_0,m_0$ and prove admissibility. At that
specialization, three tagged copies turn the query formulas into class
and pairing laws. Section~\ref{rel-diagonal} then uses these laws to
construct the terminal proof.

\subsection{Generic relations and probes}

\begin{definition}[Admissible queries and observational equivalence]
\label{rel-admissible}
Define
\[
\begin{aligned}
 P\approx P'&=\forall y.\,(P@y\leftrightarrow P'@y),\\
 i\approx i'&=\forall x.\,(i@x\leftrightarrow i'@x),\\
 \operatorname{Adm}
   &=\neg D(m)\land
       \bigwedge_{t\ {\rm structural}}\forall P,P'.\,
       \bigl(P\approx P'\Rightarrow
          (\operatorname{qry}_t(m,P)\leftrightarrow
           \operatorname{qry}_t(m,P'))\bigr),\\
 \operatorname{Good}(P)
   &=\bigl(\forall y.\,(P@y\leftrightarrow
                     P@\operatorname{Round}y)\bigr)
           \land\operatorname{qry}_G(m,P),\\
 y\sim z&=\forall P.\,
       \bigl(\operatorname{Good}(P)\Rightarrow
                      (P@y\leftrightarrow P@z)\bigr).
\end{aligned}
\]
The finite Boolean operations have fixed bracketing.  From
\(ad:\operatorname{Adm}\) obtain \(d_m:\neg D(m)\) by projection.
For each structural \(t\), projection followed by two instantiations
and application defines
\[
 \operatorname{qr}_t(e):
 \operatorname{qry}_t(m,P)\leftrightarrow
 \operatorname{qry}_t(m,P')
 \qquad(e:P\approx P').
\]
\end{definition}

\begin{lemma}[Equivalence and rounding]
\label{rel-equivalence}
The relation \(\sim\) has reflexivity, symmetry, and transitivity
builders.  There is a builder
\(\operatorname{RoundRel}(y):y\sim\operatorname{Round}y\).
For \(e:a\sim a'\) and \(f:c\sim c'\) there is a builder
\[
 \operatorname{cong}(e,f):
       (a\sim c)\leftrightarrow(a'\sim c').
\]
These builders require no admissibility assumption.
\end{lemma}
\begin{proof}
Reflexivity is
\(\Lambda P.\lambda^*h.\mathsf{pl}_{P@y\leftrightarrow P@y}()\).
For symmetry, under \(e:y\sim z\), introduce \(P,h\), instantiate
\(e[P]\cdot h\), and apply \(\mathsf{pl}\) to reverse its
biconditional.  For transitivity, under
\(e:y\sim z\) and \(f:z\sim w\), introduce \(P,h\) and apply
\(\mathsf{pl}_{P@y\leftrightarrow P@w}\) to
\(e[P]\cdot h,f[P]\cdot h\).
For rounding introduce \(P,h\), project
\(\forall w.(P@w\leftrightarrow P@\operatorname{Round}w)\)
from \(h:\operatorname{Good}(P)\), and instantiate it at \(y\).
Write these builders as \(\operatorname{refl}\),
\(\operatorname{sym}\), and \(\operatorname{trans}\).
The two components of the last biconditional are
\[
\begin{aligned}
 &\lambda^*u.\operatorname{trans}
       (\operatorname{sym}(e),\operatorname{trans}(u,f)),\\
 &\lambda^*u.\operatorname{trans}
       (e,\operatorname{trans}(u,\operatorname{sym}(f))).
\end{aligned}
\]
In the first line \(u:a\sim c\), and in the second
\(u:a'\sim c'\), so the endpoints are respectively
\(a'\sim c'\) and \(a\sim c\).
\end{proof}

\begin{definition}[Class, union, and pairing predicates]
\label{rel-pair-definition}
Put
\[
\begin{aligned}
 c_y&=\operatorname{pred}(w.w\sim y),\\
 u_{yz}&=\operatorname{pred}(w.(w\sim y)\lor(w\sim z)),\\
 \operatorname{Is}_t(y)&=\operatorname{qry}_{C_t}(m,c_y)
                         &&(t=0,\ell,r),\\
 T_t(y,z)&=\operatorname{Is}_0(y)\land
             \operatorname{Is}_t(z)\land
             \operatorname{qry}_{E_t}(m,u_{yz})
                         &&(t=\ell,r),\\
 \operatorname{form}(P,a,z)
       &=\bigl(\exists y.\,(T_\ell(y,z)\land P@y)\bigr)
                         \lor T_r(a,z),\\
 \operatorname{Pair}(P,a)&=
             \operatorname{pred}(z.\operatorname{form}(P,a,z)).
\end{aligned}
\]
These are data definitions, independent of any proof of
\(\operatorname{Adm}\).
\end{definition}

\begin{lemma}[Congruence of the relational predicates]
\label{rel-pair-congruence}
There are unrounded evaluation proofs
\[
 c_y@w\leftrightarrow(w\sim y),\qquad
 u_{yz}@w\leftrightarrow((w\sim y)\lor(w\sim z)).
\]
Equivalent parameters give \(c_y\approx c_{y'}\) and
\(u_{yz}\approx u_{y'z'}\).  Under \(ad:\operatorname{Adm}\),
the predicates \(\operatorname{Is}_t,T_t\) respect \(\sim\).
Moreover, still under \(ad\), and under \(P\approx P'\),
\(a\sim a'\), and \(z\sim z'\)
there are proofs of
\[
 \operatorname{form}(P,a,z)\leftrightarrow
       \operatorname{form}(P',a',z'),
\]
and, under the first two of these assumptions, of
\(\operatorname{Pair}(P,a)\approx\operatorname{Pair}(P',a')\).
In particular, under \(ad\) there is
\[
 \operatorname{pv}(P,a,z):
  \operatorname{Pair}(P,a)@z
       \leftrightarrow\operatorname{form}(P,a,z).
\]
\end{lemma}
\begin{proof}
Predicate evaluation first gives
\[
 c_y@w\leftrightarrow(\operatorname{Round}w\sim y).
\]
Apply relation congruence with
\(\operatorname{RoundRel}(w)\) and \(\operatorname{refl}(y)\);
propositional chaining gives the first displayed evaluation.
For the union do this separately for \(y\) and \(z\), then chain
their disjunctions.  For \(e:y\sim y'\), congruence with
\(\operatorname{refl}(w),e\), together with the two class
evaluations, gives \(c_y@w\leftrightarrow c_{y'}@w\).
Introduce \(w\) to obtain the class comparison.  For the union
use the two parameter congruences inside the disjunction before
introducing \(w\).

Apply \(\operatorname{qr}_{C_t}\) to the class comparison to
obtain the \(\operatorname{Is}_t\) comparison.  For \(T_t\),
combine the two resulting \(\operatorname{Is}\) comparisons and
\(\operatorname{qr}_{E_t}\) applied to the union comparison.
For \(\operatorname{form}\), fix \(y\).  The \(T_\ell\)
comparison with its first argument unchanged, together with
\((P\approx P')[y]\), gives
\[
 (T_\ell(y,z)\land P@y)\leftrightarrow
           (T_\ell(y,z')\land P'@y).
\]
Apply \(\mathsf{Ex}_y\), then combine its result with the
\(T_r(a,z)\leftrightarrow T_r(a',z')\) comparison by
\(\mathsf{pl}\).  Predicate evaluation for \(\operatorname{Pair}\)
has right side
\(\operatorname{form}(P,a,\operatorname{Round}z)\).
The just-proved form congruence, with \(P,a\) unchanged and
\(\operatorname{RoundRel}(z)\), gives \(\operatorname{pv}\).
Finally fix \(z\), combine both \(\operatorname{pv}\) proofs
with form congruence at unchanged \(z\), and introduce \(z\).
\end{proof}

\begin{definition}[Lifts and uniform probes]
\label{rel-lifts}
Define data, still with open \(A,m\), by
\[
\begin{aligned}
 \operatorname{Part}(v)&=\operatorname{pred}(y.v\diamond_m y),\\
 \operatorname{sb}(v,m,a)&=\operatorname{qry}_R
                    (m,\operatorname{Pair}(\operatorname{Part}(v),a)),\\
 \operatorname{lift}_m(a)&=\operatorname{mk}_X(v.\operatorname{sb}(v,m,a)).
\end{aligned}
\]
For \(i:\mathsf P_X\), define \(z_i:V\) by the following data
lambda.  At \(z:Z\), pack along the even path \(g\to J\) the
data function
\[
 \lambda^{\mathsf d,\mathsf d}w:
             \operatorname{Single}(z^\downarrow).
 \operatorname{send}_n
    \bigl(\operatorname{Access}
          (w;o,y.i@\operatorname{lift}_o(y))\bigr),
\]
where all channel types and the lift in its body use
\(A=z^\downarrow\).  This defines the body of
\(\lambda^{\mathsf d,\mathsf d}z:Z.\,\cdots\).
The variable \(o\) is the data parameter supplied by
\(\operatorname{Access}\).  Thus \(z_i\) has no free \(A,m\).
Put
\[
\begin{aligned}
 \operatorname{le}(i,x)&=[x z_i],\\
 x\simeq x'&=\forall i.\,(\operatorname{le}(i,x)\leftrightarrow \operatorname{le}(i,x')),\\
 \operatorname{Valid}(x)&=\forall i,i'.\,
       \bigl(i\approx i'\Rightarrow
                    (\operatorname{le}(i,x)\leftrightarrow \operatorname{le}(i',x))\bigr),\\
 \operatorname{Ext}(i)&=\forall x,x'.\,
       \bigl(x\simeq x'\Rightarrow(i@x\leftrightarrow i@x')\bigr),\\
 \operatorname{Resp}(v)&=\forall y,z.\,
       \bigl(y\sim z\Rightarrow
                    (v\diamond_m y\leftrightarrow v\diamond_m z)\bigr).
\end{aligned}
\]
\end{definition}

\begin{lemma}[Evaluation, validity, and response]
\label{rel-lift-properties}
The relation \(\simeq\) is an equivalence relation, with the
congruence builder of Lemma~\ref{rel-equivalence}, and
\[
 \operatorname{RoundRel}_X(x):x\simeq\operatorname{Round}_Xx.
\]
Data cancellation supplies
\[
 \operatorname{LP}(i,a):
       \operatorname{le}(i,\operatorname{lift}_m a)\leftrightarrow \operatorname{sb}(z_i,m,a).
\]
Under \(ad:\operatorname{Adm}\) there are builders
\begin{align}
 z_i\diamond_m y&\leftrightarrow
             i@\operatorname{lift}_m(\operatorname{Round}y),
                  \label{rel-probe-rounded}\\
 \operatorname{LC}(e)&:
       \operatorname{lift}_m(a)\simeq\operatorname{lift}_m(a')
                    \qquad(e:a\sim a'), \label{rel-lift-congruence}\\
 &\operatorname{Valid}(\operatorname{lift}_m a). \label{rel-lift-valid}
\end{align}
If also \(ex:\operatorname{Ext}(i)\), then
\begin{equation}
 z_i\diamond_m y\leftrightarrow i@\operatorname{lift}_m y
 \quad\hbox{and}\quad \operatorname{Resp}(z_i).
 \label{rel-probe-exact}
\end{equation}
If \(h:\operatorname{Resp}(v)\), then
\begin{equation}
 \operatorname{Part}(v)@y\leftrightarrow v\diamond_m y.
 \label{rel-part-exact}
\end{equation}
\end{lemma}
\begin{proof}
For \(\simeq\), introduce \(i\) and use propositional
reflexivity, symmetry, or transitivity of the instantiated
biconditionals.  Its congruence builder is the pair displayed
in Lemma~\ref{rel-equivalence}, with the relation replaced by
\(\simeq\).  Primitive \(X\)-evaluation is preserved by
\(\operatorname{Round}_X\).  Apply this preservation at \(z_i\)
and introduce \(i\) to obtain \(\operatorname{RoundRel}_X\).
Cancellation of \(\operatorname{mk}_X\) gives \(\operatorname{LP}\)
at the exact input \(z_i\).

For \eqref{rel-probe-rounded}, the application \(z_iA^\uparrow\)
beta-reduces to the same pack and data lambda at \(A\), since
\((A^\uparrow)^\downarrow=_\beta A\) and the wrapper choices
are fixed.  The pack/use cancellation leaves the read of that
function applied to \(\operatorname{Bundle}(m,y)\).
After its raw beta contraction, read/send cancellation leaves
\[
 \operatorname{Access}
   (\operatorname{Bundle}(m,y);o,w.i@\operatorname{lift}_o(w)).
\]
The Access builder, applied to \(d_m\), compares this with
\(i@\operatorname{lift}_m(\operatorname{Round}y)\).
Chain these three biconditionals by \(\mathsf{pl}\).

To obtain \(\operatorname{LC}(e)\), introduce \(i\).
Pair congruence applied to the reflexive comparison of
\(\operatorname{Part}(z_i)\) and to \(e\), followed by
\(\operatorname{qr}_R\), compares \(\operatorname{sb}(z_i,m,a)\) and
\(\operatorname{sb}(z_i,m,a')\).  Combine it with the two \(\operatorname{LP}\)
proofs and introduce \(i\).

For validity introduce \(i,i'\) and \(e:i\approx i'\).
There is a comparison
\(\operatorname{Part}(z_i)\approx\operatorname{Part}(z_{i'})\):
at \(y\), predicate evaluation gives the two expressions
\(z_i\diamond_m\operatorname{Round}y\) and
\(z_{i'}\diamond_m\operatorname{Round}y\).
Apply \eqref{rel-probe-rounded} at precisely
\(\operatorname{Round}y\), and use
\[
 e[\operatorname{lift}_m
               (\operatorname{Round}(\operatorname{Round}y))].
\]
These proofs propositionally entail the required evaluation
comparison; introduce \(y\).  Apply Pair congruence with \(a\)
unchanged, then \(\operatorname{qr}_R\), and finally the two
\(\operatorname{LP}\) proofs.  Introducing \(i,i',e\) proves
\eqref{rel-lift-valid}.

For \eqref{rel-probe-exact}, apply \(ex\) at
\(\operatorname{lift}_m y,\operatorname{lift}_m(\operatorname{Round}y)\)
to \(\operatorname{LC}(\operatorname{RoundRel}(y))\).
Combine its result with \eqref{rel-probe-rounded}.  Under
\(e:y\sim z\), apply \(ex\) at the two lifts to
\(\operatorname{LC}(e)\), and combine this with the exact
probe evaluations at \(y,z\).  Introducing \(y,z,e\) proves
\(\operatorname{Resp}(z_i)\).
Finally predicate evaluation gives
\(\operatorname{Part}(v)@y\leftrightarrow
v\diamond_m\operatorname{Round}y\).
Combine it with \(h[y][\operatorname{Round}y]\cdot
\operatorname{RoundRel}(y)\) to obtain \eqref{rel-part-exact}.
\end{proof}

\subsection{The laws needed for the diagonal argument}
\label{rel-target-laws}

The relations and uniform probes are now defined.  Before choosing
$A$ and $m$, we identify the laws their specialization must provide.
For $i:\mathsf P_X$, define
\begin{equation}
 \operatorname{Ind}(i)=
 \forall x:X.\bigl(\operatorname{Valid}(x)\Rightarrow
       (\operatorname{le}(i,x)\Rightarrow i@x)\bigr).
 \label{rel-ind-definition}
\end{equation}
This formula says that $i$ contains every valid $x$ for which
$\operatorname{le}(i,x)$ holds.  The name records its role in the
well-foundedness diagonal; it does not assume a pre-existing order on $X$.

We shall construct data operations $x\mapsto\delta x:X$ and
$i\mapsto Si:\mathsf P_X$, and a data term
$\operatorname{WF}:X$.  Their required properties are
\begin{align*}
 &\operatorname{Valid}(\delta x),\qquad
                    \operatorname{Valid}(\operatorname{WF}),\\
 &\operatorname{Ext}(i)\ \Longrightarrow\
     \operatorname{Ext}(Si)\ \text{ and }\
                  (Si@x\leftrightarrow i@\delta x),\\
 &\operatorname{Ext}(i),\ \operatorname{Valid}(x)
       \ \Longrightarrow\
       \bigl(\operatorname{le}(i,\delta x)
                      \leftrightarrow\operatorname{le}(Si,x)\bigr),\\
 &\operatorname{Ext}(i)\ \Longrightarrow\
       \bigl(\operatorname{le}(i,\operatorname{WF})
                              \leftrightarrow\operatorname{Ind}(Si)\bigr).
\end{align*}
Here the outer implication signs describe proof builders with the
displayed inputs, as elsewhere in the construction.  We also require
$x\simeq x'$ to imply $\delta x\simeq\delta x'$.
The logical binders introduced in the final diagonal proof range only
over $X$ or $\mathsf P_X$; it also uses the previously proved laws at
their original domains.

These laws already explain the intended role of $\operatorname{WF}$.
For an extensional $i$, a proof of $\operatorname{Ind}(i)$ gives a
proof of $\operatorname{Ind}(Si)$: apply the proof of
$\operatorname{Ind}(i)$ at
the valid point $\delta x$, using the predecessor comparison to turn
$\operatorname{le}(Si,x)$ into $\operatorname{le}(i,\delta x)$.
The shift evaluation then gives the required conclusion $Si@x$.
The last displayed law then gives
$\operatorname{le}(i,\operatorname{WF})$, so validity of
$\operatorname{WF}$ gives $i@\operatorname{WF}$.  Section~\ref{rel-diagonal}
uses this consequence with a specified negative predicate to derive
$\bot$; it will prove the construction with every guard in place.

To obtain the predecessor comparison, we will encode three distinguishable
copies of $X$ in $L$.  The class and edge queries will identify the
corresponding elements of those copies.  Observations of the pairing
predicate on one copy will be equivalent to the evaluations of the
represented predicate; observations on another will identify the point
of evaluation up to $\simeq$.  The $R$ query can then express
$\operatorname{le}$ at that point.  Its validity guard is essential:
it permits transport between the corresponding
$\operatorname{le}$ formulas for the represented and original predicates.
The next two subsections build
these copies and queries before the guarded diagonal uses them.

\subsection{Three tagged copies and one admissible parameter}

The types $V,X,\mathsf P_X$ and the probes $z_i$ were constructed
uniformly, before any choice of $A,m$.  We now use them to choose
$A_0$ and then $m_0$.  The observations defining $m_0$ may use those
already constructed probes; they never use $m_0$ itself.  This order
will allow us to prove admissibility after defining the data term.

We now choose a fixed \(A_0:s\).  Choose the secondary path
\(k\to d\), inside \(C\), so that its concatenation with the
previously fixed \(d\to k\) has even parity and at least four
negative edges, and put \(A_0=W_{k\to d}(V)\).
For completeness, in a component with an odd closed walk a
detour changes the parity if necessary; in a component with a
consistent two-coloring the return paths already have matching
parities.  A closed detour through an internal negative exists
by strong connectivity.  Inserting it twice preserves parity
and increases the number of negatives.  Repeating this insertion
gives the claimed choice.

Split the resulting wrapper from \(V\) to \(B(A_0)\) into
four odd segments \(\rho_1,\ldots,\rho_4\).  This is possible
by cutting after the first, second, and third negative;
the remaining number of negatives is odd.  Let \(T_0=V\) and
\(T_a=W_{\rho_a}(T_{a-1})\), so \(T_4=B(A_0)\).
The fixed odd path \(k\to j\), denoted \(\rho_5\) here, gives
\(T_5=L(A_0)\).  The first odd wrapper defines
\[
 y_1(x)=\operatorname{mk}_{\rho_1}(v.[xv]):T_1.
\]
Use the double-embedding construction of Definition~\ref{chan-double}
and Lemma~\ref{chan-double-laws} on \((\rho_2,\rho_3)\)
and on \((\rho_4,\rho_5)\).  Denote their signed embeddings
by \(E^{(1)}_\epsilon:T_1\to T_3\) and
\(E^{(2)}_\epsilon:T_3\to T_5\), their bit observations by
\(B_1,B_2\), and their recovery formula builders by
\(\mathcal R_1,\mathcal R_2\).  These are notation for the
previous data macros, not logical function spaces.
Fix distinct strings, for example
\[
 (b_1(0),b_2(0))=(0,0),\quad
 (b_1(\ell),b_2(\ell))=(1,0),\quad
 (b_1(r),b_2(r))=(0,1),
\]
and define
\[
\begin{aligned}
 x^t&=E^{(2)}_{b_2(t)}(E^{(1)}_{b_1(t)}(y_1(x))),\\
 H_2(y)&=B_2(y),\\
 H_1(y)&=\mathcal R_2(y;c.B_1(c)),\\
 O_i(y)&=\mathcal R_2
       (y;c.\mathcal R_1(c;u.\langle u,z_i\rangle_{\rho_1})).
\end{aligned}
\]
Here \(\langle u,z_i\rangle_{\rho_1}\) uses the first odd
wrapper, whose input type is \(V\).

\begin{lemma}[Tagged observations]
\label{rel-tags}
For \(t=0,\ell,r\) there are proofs of the signed bits
\(H_a(x^t)\) if \(b_a(t)=1\), and \(\neg H_a(x^t)\) if
\(b_a(t)=0\), and proofs
\[
 O_i(x^t)\leftrightarrow \operatorname{le}(i,x).
\]
Each \(H_a\) and \(O_i\) is invariant, by biconditional, under
\(y\mapsto\operatorname{Round}y\).
Define
\[
 \operatorname{Ae}(y,z)=
   \bigwedge_{a=1,2}(H_a(y)\leftrightarrow H_a(z))
    \land\bigl(\forall i.\,(O_i(y)\leftrightarrow O_i(z))\bigr).
\]
Then \(\operatorname{Ae}\) is an equivalence relation and
there is \(\operatorname{AeRound}(y):
\operatorname{Ae}(y,\operatorname{Round}y)\).
For each tag \(t\), its restriction compares exactly \(\simeq\):
\[
 \operatorname{Ae}(x^t,x'^{\,t})\leftrightarrow x\simeq x'.
\]
For distinct tags \(t,u\) there is a proof of
\(\neg\operatorname{Ae}(x^t,x'^{\,u})\).
\end{lemma}
\begin{proof}
The two double-embedding builders give their signed bits and
recover each callback at its exact input.  Apply them successively
to \(B_1\) and to \(u\mapsto\langle u,z_i\rangle_{\rho_1}\).
For the latter, first-wrapper cancellation gives
\(\langle y_1(x),z_i\rangle_{\rho_1}\leftrightarrow[xz_i]\).
Propositional chaining gives all the asserted canonical
evaluations.  For a general \(y\), the outer bit and every
outer recovery are Boolean combinations of primitive
\(L\)-evaluations at arguments independent of \(y\).
Primitive Round preservation at those arguments, followed
by \(\mathsf{pl}\), therefore gives each invariance proof.

Reflexivity of \(\operatorname{Ae}\) pairs the propositional
bit reflexivities with
\(\Lambda i.\mathsf{pl}_{O_i(y)\leftrightarrow O_i(y)}()\).
For symmetry, project both bit comparisons and the observation
universal, reverse the former by \(\mathsf{pl}\), instantiate
the latter at \(i\), reverse it, and introduce \(i\).
For transitivity perform these projections on both premises,
compose the bit comparisons, and at each \(i\) compose the
two observation comparisons before introducing \(i\).
Reassemble the conjunctions by \(\mathsf{pl}\).
The invariance proofs supply the two bit components of
\(\operatorname{AeRound}\) and, after introduction of \(i\),
its observation component.  The congruence builder for
\(\operatorname{Ae}\) is the explicit two-component
transitivity construction of Lemma~\ref{rel-equivalence}.

Given \(e:x\simeq x'\), the canonical signed bits imply
the two same-tag bit comparisons.  At \(i\), combine \(e[i]\)
with both canonical \(O_i\) evaluations, then introduce \(i\).
This proves \(\operatorname{Ae}(x^t,x'^{\,t})\).
Conversely project its observation universal, instantiate at
\(i\), and chain the two canonical evaluations to obtain
\(\operatorname{le}(i,x)\leftrightarrow \operatorname{le}(i,x')\); introduce \(i\).
For \(t\ne u\), choose a coordinate where their fixed strings
differ.  The two signed-bit proofs contradict the comparison
at that coordinate projected from an assumed
\(\operatorname{Ae}(x^t,x'^{\,u})\).
Abstract that assumption to obtain the negation.
\end{proof}

\begin{definition}[The specialized query object]
\label{rel-specialized-object}
At \(A_0\), for \(P:\mathsf P_L\), define
\[
\begin{aligned}
 \mathcal P_G(P)&=\forall y,z.\,
  \bigl(\operatorname{Ae}(y,z)\Rightarrow(P@y\leftrightarrow P@z)\bigr),\\
 \mathcal P_{C_t}(P)&=\exists x.\forall w.\,
                    (P@w\leftrightarrow\operatorname{Ae}(w,x^t)),\\
 \mathcal P_{E_t}(P)&=\exists x.\forall w.\,
       \bigl(P@w\leftrightarrow
        (\operatorname{Ae}(w,x^0)\lor\operatorname{Ae}(w,x^t))\bigr),\\
 i_P&=\operatorname{pred}_X(x.P@(x^\ell)),\\
 \mathcal P_R(P)&=\operatorname{Ext}(i_P)\Rightarrow
     \forall x.\bigl(\operatorname{Valid}(x)\Rightarrow
                    (P@(x^r)\Rightarrow \operatorname{le}(i_P,x))\bigr).
\end{aligned}
\]
Here \(C_t\) has \(t=0,\ell,r\) and \(E_t\) has \(t=\ell,r\).
Define
\[
 m_0=\operatorname{mk}_M
   \left(bp.\bigvee_{t\ {\rm structural}}
       \bigl(\operatorname{Mark}_t(bp)\land
                       \mathcal P_t(\operatorname{Dec}(bp))\bigr)\right).
\]
The formulas \(\mathcal P_t\) use \(z_i\), tags, and observations
already defined independently of \(m_0\); this is a data term
with no recursive occurrence of itself.
\end{definition}

\begin{lemma}[Respect and admissibility]
\label{rel-query-realization}
For \(d:i\approx i'\) there is a builder
\[
 \operatorname{ExtC}(d):
       \operatorname{Ext}(i)\leftrightarrow\operatorname{Ext}(i').
\]
For \(e:P\approx P'\), every structural \(t\) has a builder
\(\mathcal P_t(P)\leftrightarrow\mathcal P_t(P')\).
Consequently there are proofs
\[
 q_t(P):\operatorname{qry}_t(m_0,P)\leftrightarrow\mathcal P_t(P)
 \quad(t\ {\rm structural}),\qquad
 ad:\operatorname{Adm}\big|_{A=A_0,m=m_0}.
\]
\end{lemma}
\begin{proof}
Under \(x,x'\), the evaluations \(d[x],d[x']\) propositionally
entail
\[
 (x\simeq x'\Rightarrow(i@x\leftrightarrow i@x'))
 \leftrightarrow
 (x\simeq x'\Rightarrow(i'@x\leftrightarrow i'@x')).
\]
Apply \(\mathsf{All}_{x,x'}\) to obtain
\(\operatorname{ExtC}(d)\).

For \(G\), under \(y,z\), apply \(\mathsf{pl}\) to
\(e[y],e[z]\) with target the comparison of the two implications
whose common antecedent is \(\operatorname{Ae}(y,z)\);
then apply \(\mathsf{All}_{y,z}\).
For \(C_t\), under \(x,w\), the proof \(e[w]\) compares
\[
 P@w\leftrightarrow\operatorname{Ae}(w,x^t)
 \quad\hbox{with}\quad
 P'@w\leftrightarrow\operatorname{Ae}(w,x^t).
\]
Apply \(\mathsf{All}_w\) and then \(\mathsf{Ex}_x\).
For \(E_t\), use exactly the same two typed operations, with
their common right side now
\(\operatorname{Ae}(w,x^0)\lor\operatorname{Ae}(w,x^t)\).

For \(R\), predicate evaluation at \(x\) for \(i_P,i_{P'}\)
has right sides
\[
 P@((\operatorname{Round}_Xx)^\ell)
 \quad\hbox{and}\quad P'@((\operatorname{Round}_Xx)^\ell).
\]
Combine these evaluations with
\(e[(\operatorname{Round}_Xx)^\ell]\) and introduce \(x\),
obtaining \(d:i_P\approx i_{P'}\).
The proof \(\operatorname{ExtC}(d)\) compares the two
antecedents of \(\mathcal P_R\).
Under \(x,v:\operatorname{Valid}(x)\), combine
\[
 e[x^r],\qquad v[i_P][i_{P'}]\cdot d
\]
with target
\[
 (P@(x^r)\Rightarrow \operatorname{le}(i_P,x))
   \leftrightarrow(P'@(x^r)\Rightarrow \operatorname{le}(i_{P'},x)).
\]
Apply \(\mathsf{Gd}_{v:\operatorname{Valid}(x)}\), then
\(\mathsf{All}_x\).  This compares the consequents of
\(\mathcal P_R\).  Combine it propositionally with
\(\operatorname{ExtC}(d)\).  In particular, the change of
predicate in \(le\) has used its explicit validity guard.

Data cancellation at \(\operatorname{Enc}_t(P)\) compares
\(\operatorname{qry}_t(m_0,P)\) with the defining disjunction
evaluated there.  The signed marker proofs eliminate every
disjunct except \(t\), and give a comparison with
\(\mathcal P_t(\operatorname{Dec}(\operatorname{Enc}_t(P)))\).
Introduce \(y\) in the decoder evaluation comparison to get
\[
 \operatorname{Dec}(\operatorname{Enc}_t(P))\approx P.
\]
Use the property-respect builder just proved and chain to
obtain \(q_t(P)\).  Cancellation at \(bp_\bot\) and the
negations of all its structural markers give \(\neg D(m_0)\).
For each structural \(t\), form
\[
 \Lambda P,P'.\lambda^*e.\,
 \mathsf{pl}\bigl(q_t(P),q_t(P'),
                      \text{property-respect}_t(e)\bigr)
\]
with target
\(\operatorname{qry}_t(m_0,P)\leftrightarrow
\operatorname{qry}_t(m_0,P')\).
Together with \(\neg D(m_0)\), these are exactly the
conjuncts of the displayed admissibility formula.
\end{proof}

From this point onward \(A=A_0\), \(m=m_0\), and \(ad\) is
the proof just constructed.

\subsection{Pairing laws at the fixed specialization}

\begin{lemma}[The coded equivalence is the observation equivalence]
\label{rel-observation-equivalence}
There are proofs
\[
 \operatorname{GA}(P):
       \operatorname{Good}(P)\leftrightarrow\mathcal P_G(P),
 \qquad
 \operatorname{SA}(y,z):
       (y\sim z)\leftrightarrow\operatorname{Ae}(y,z).
\]
For each \(t\), there are builders
\[
 \operatorname{Tag}_t(e):x^t\sim x'^{\,t}
                  \quad(e:x\simeq x'),\qquad
 \operatorname{TagBack}_t(f):x\simeq x'
                  \quad(f:x^t\sim x'^{\,t}).
\]
For \(t\ne u\) there is
\(\operatorname{diff}_{t,u}:\neg(x^t\sim x'^{\,u})\).
\end{lemma}
\begin{proof}
The proof \(q_G(P)\) identifies the second conjunct of
\(\operatorname{Good}(P)\).  Its property implies the first:
the builder
\[
 \lambda^*g.\Lambda y.\,
      g[y][\operatorname{Round}y]\cdot\operatorname{AeRound}(y)
\]
has that implication as type.  Combining these by
\(\mathsf{pl}\) gives \(\operatorname{GA}(P)\).

For an observation \(\Phi\) equal to \(H_1,H_2\), or \(O_i\)
with \(i\) temporarily free, put
\(P_\Phi=\operatorname{pred}(w.\Phi(w))\).
Predicate evaluation and Round invariance give, at every \(u\),
an evaluation proof \(P_\Phi@u\leftrightarrow\Phi(u)\).
Under \(a:\operatorname{Ae}(u,w)\), project the relevant bit
comparison, or project and instantiate its observation
universal at the free \(i\).  Combine it with the two evaluations
to get \(P_\Phi@u\leftrightarrow P_\Phi@w\).
Introducing \(u,w,a\) proves \(\mathcal P_G(P_\Phi)\);
apply \(\operatorname{back}(\operatorname{GA}(P_\Phi),-)\)
to get \(\operatorname{Good}(P_\Phi)\).

Under \(e:y\sim z\), the proof
\(e[P_\Phi]\cdot\operatorname{Good}(P_\Phi)\), combined
with the two evaluations, gives \(\Phi(y)\leftrightarrow\Phi(z)\).
Take the two bit instances.  For \(O_i\), introduce \(i\)
after obtaining this comparison.  Their conjunction proves
\(\operatorname{Ae}(y,z)\).  Conversely, under
\(a:\operatorname{Ae}(y,z)\), the term
\[
 \Lambda P.\lambda^*g.\,
       \bigl(\operatorname{to}(\operatorname{GA}(P),g)\bigr)
                       [y][z]\cdot a
\]
proves \(y\sim z\).  Abstract the two premises and pair the
directions to obtain \(\operatorname{SA}\).

Given \(e:x\simeq x'\), use the same-tag direction of
Lemma~\ref{rel-tags}, then \(\operatorname{back}(\operatorname{SA},-)\).
Given \(f:x^t\sim x'^{\,t}\), use
\(\operatorname{to}(\operatorname{SA},f)\) and the reverse
same-tag direction.  For distinct tags, introduce
\(f:x^t\sim x'^{\,u}\), map it forward through
\(\operatorname{SA}\), and apply the negation from
Lemma~\ref{rel-tags}.  These define the three asserted builders.
\end{proof}

\begin{lemma}[Class and edge witnesses]
\label{rel-witnesses}
For \(t=0,\ell,r\) there is a proof
\[
 \operatorname{Is}_t(y)\leftrightarrow
                         \exists x.\,(y\sim x^t).
\]
For \(t=\ell,r\) there is a proof
\[
 T_t(y,z)\leftrightarrow
             \exists x.\,((y\sim x^0)\land(z\sim x^t)).
\]
\end{lemma}
\begin{proof}
For the first forward implication, take
\(h:\operatorname{Is}_t(y)\), apply
\(\operatorname{to}(q_{C_t}(c_y),h)\), and open the result
as \(x,p\), where
\[
 p:\forall w.\,(c_y@w\leftrightarrow\operatorname{Ae}(w,x^t)).
\]
The class evaluation and reflexivity prove \(c_y@y\).
Apply \(p[y]\) forward and \(\operatorname{SA}(y,x^t)\)
backward to obtain \(y\sim x^t\).  Witness the opened \(x\)
in the target existential.  The target contains no free
occurrence of either local witness.
For the reverse, open \(x,e:y\sim x^t\).
At \(w\), the class evaluation, congruence with
\(\operatorname{refl}(w),e\), and
\(\operatorname{SA}(w,x^t)\) yield
\(c_y@w\leftrightarrow\operatorname{Ae}(w,x^t)\).
Introduce \(w\), witness \(x\) in \(\mathcal P_{C_t}(c_y)\),
and apply \(q_{C_t}(c_y)\) backward.

For the edge forward implication, take \(h:T_t(y,z)\).
Project both \(\operatorname{Is}\) conjuncts and its query.
Use the first equivalence and open their witnesses as
\[
 a,e_a:y\sim a^0,\qquad
 b',e_{b'}:z\sim b'^{\,t}.
\]
Map the query forward through \(q_{E_t}(u_{yz})\) and
open its witness as \(x,p\), with
\[
 p:\forall w.\,
   \bigl(u_{yz}@w\leftrightarrow
    (\operatorname{Ae}(w,x^0)\lor\operatorname{Ae}(w,x^t))\bigr).
\]
The union evaluation and reflexivity prove \(u_{yz}@y\)
and \(u_{yz}@z\).  Apply \(p\) at each of these points and
the comparisons \(\operatorname{SA}\); the resulting proofs
have types
\[
 (y\sim x^0)\lor(y\sim x^t),\qquad
 (z\sim x^0)\lor(z\sim x^t).
\]
The term
\[
 \lambda^*f.\operatorname{diff}_{0,t}\cdot
        \operatorname{trans}(\operatorname{sym}(e_a),f)
\]
negates \(y\sim x^t\).  The corresponding term
\[
 \lambda^*f.\operatorname{diff}_{t,0}\cdot
        \operatorname{trans}(\operatorname{sym}(e_{b'}),f)
\]
negates \(z\sim x^0\).  Propositional elimination gives
\(y\sim x^0\) and \(z\sim x^t\); pair them and witness \(x\).
All three openings have the original target existential.

For the reverse, open \(x,e\) and project
\(e_0:y\sim x^0\), \(e_t:z\sim x^t\).
Witness \(x\) with \(e_0\) or \(e_t\) in the first
equivalence backward to obtain the two \(\operatorname{Is}\)
conjuncts.  At \(w\), the union evaluation and congruences
along \(e_0,e_t\) compare \(u_{yz}@w\) with
\((w\sim x^0)\lor(w\sim x^t)\).  Use the two
\(\operatorname{SA}\) proofs and introduce \(w\), obtaining
the property required for the witness \(x\) in
\(\mathcal P_{E_t}(u_{yz})\).  Apply \(q_{E_t}(u_{yz})\)
backward and conjoin the three results.
\end{proof}

\begin{lemma}[The two pairing laws]
\label{rel-pair-laws}
Under \(h:\operatorname{Resp}(v)\), let
\(P_x=\operatorname{Pair}(\operatorname{Part}(v),x^0)\).
For every \(x':X\) there are proofs
\begin{align}
 P_x@(x'^{\,\ell})&\leftrightarrow
                         v\diamond_m(x'^{\,0}), \label{rel-pair-left}\\
 P_x@(x'^{\,r})&\leftrightarrow x\simeq x'. \label{rel-pair-right}
\end{align}
\end{lemma}
\begin{proof}
For \eqref{rel-pair-left} forward, first negate
\(T_r(x^0,x'^{\,\ell})\).  Given such a proof, map it
forward through Lemma~\ref{rel-witnesses} and open its
witness \(u\) into \(\bot\).  Its second conjunct is
\(x'^{\,\ell}\sim u^r\), contradicted by
\(\operatorname{diff}_{\ell,r}\).
Given \(p:P_x@(x'^{\,\ell})\), use
\(\operatorname{pv}\) and this negation to extract
\[
 \exists y.\,
       (T_\ell(y,x'^{\,\ell})\land\operatorname{Part}(v)@y).
\]
Open it as \(y,q\), and open the edge witness from the
first conjunct as \(u,q_0,q_\ell\), where
\[
 q_0:y\sim u^0,\qquad q_\ell:x'^{\,\ell}\sim u^\ell.
\]
The proof \(\operatorname{TagBack}_\ell(q_\ell)\) is
\(x'\simeq u\).  Symmetrize it and apply
\(\operatorname{Tag}_0\), then compose with \(q_0\),
to get \(e:y\sim x'^{\,0}\).
From the second conjunct of \(q\), the unrounded Part
evaluation gives \(v\diamond_m y\).
Apply \(h[y][x'^{\,0}]\cdot e\) forward to obtain the
target \(v\diamond_m(x'^{\,0})\).  Each opening uses this
fixed target.

For the left reverse implication, a proof of
\(v\diamond_m(x'^{\,0})\) gives
\(\operatorname{Part}(v)@(x'^{\,0})\) by
\eqref{rel-part-exact} backward.  Witness \(x'\) with
the two reflexivities in the edge equivalence backward
to get \(T_\ell(x'^{\,0},x'^{\,\ell})\).
Pair these, witness \(y=x'^{\,0}\), inject the resulting
existential into \(\operatorname{form}\), and apply
\(\operatorname{pv}\) backward.

For \eqref{rel-pair-right} forward, negate
\(\exists y.(T_\ell(y,x'^{\,r})\land\operatorname{Part}(v)@y)\).
Open an assumed existential into \(\bot\), open its
\(T_\ell\) witness, and contradict its second conjunct
\(x'^{\,r}\sim u^\ell\) with \(\operatorname{diff}_{r,\ell}\).
Thus \(\operatorname{pv}\) maps a proof of
\(P_x@(x'^{\,r})\) to \(T_r(x^0,x'^{\,r})\).
Open its witness \(u\), with
\(x^0\sim u^0\) and \(x'^{\,r}\sim u^r\).
The two same-tag converses give \(x\simeq u\) and
\(x'\simeq u\).  Compose the first with the symmetry of
the second to prove \(x\simeq x'\).
For the reverse, given \(e:x\simeq x'\), witness \(x\)
in the edge equivalence with reflexivity at \(x^0\) and
\(\operatorname{Tag}_r(\operatorname{sym}(e)):
x'^{\,r}\sim x^r\).  This gives \(T_r(x^0,x'^{\,r})\);
inject it into \(\operatorname{form}\), then apply
\(\operatorname{pv}\) backward.
Abstract each input proof and pair the two directions.
\end{proof}

\section{The guarded diagonal argument}
\label{rel-diagonal}

The pairing laws are now available at the fixed specialization. We use
them to compare a probe with a predicate on $X$, and then derive a proof
of $F_b$.

The construction below adapts the well-foundedness paradox of Hurkens
\cite{Hurkens1995}, as presented by Geuvers~\cite[Section~2]{Geuvers2007}.
We prove each identity with the validity and extensionality hypotheses
required by the present encoding.

We retain the specialization $A=A_0$, $m=m_0$ and the constructed
proof $ad:\operatorname{Adm}$.  Thus every relation and every occurrence
of $\diamond_m$ below uses this fixed specialization.  We use the Pair
laws of Lemma~\ref{rel-pair-laws}, the query comparison $q_R$, and the
previously constructed lift and tag builders.  All displayed proof
builders are in the labelled fragment.  In particular, $\Lambda$ and
$\lambda^*$ denote the established logical introduction rules, and
$\mathsf{pl}_F$ denotes the established propositional proof builder
with target formula $F$.  Its use to construct a biconditional supplies
both implication directions; it never replaces a data expression by
a logically equivalent expression.  The data parameters $v:V$ used
below are open parameters, not additional logical quantifier domains.

\begin{lemma}[The predicate represented by a response]
\label{rel-response-predicate}
For every data parameter $v:V$, define
\[
 i_v=\operatorname{pred}_X(x'.\,v\diamond_m(x'^{\,0}))
       :\mathsf P_X.
\]
This data definition requires no proof hypothesis.  Given
$h:\operatorname{Resp}(v)$, there are builders
\[
 E_h(x):i_v@x\leftrightarrow v\diamond_m(x^0),
 \qquad ex_v:\operatorname{Ext}(i_v).
\]
Moreover, for $x:X$ and $u:\operatorname{Valid}(x)$, there is a builder
\begin{equation}
 \operatorname{SB}(v,h,x,u):
 \operatorname{sb}(v,m,x^0)\leftrightarrow
                         \operatorname{le}(i_v,x).
 \label{rel-sb-law}
\end{equation}
\end{lemma}

\begin{proof}
The predicate evaluation law gives
\[
 i_v@x\leftrightarrow
      v\diamond_m((\operatorname{Round}_Xx)^0).
\]
On the other hand,
\[
 h[x^0][(\operatorname{Round}_Xx)^0]\cdot
       \operatorname{Tag}_0(\operatorname{RoundRel}_X(x))
 :v\diamond_m(x^0)\leftrightarrow
      v\diamond_m((\operatorname{Round}_Xx)^0).
\]
Applying $\mathsf{pl}$ to these two biconditionals at the stated
target gives $E_h(x)$.  For $e:x\simeq x'$, put
\[
 ex_v=\Lambda x,x'.\lambda^*e.\,
 \mathsf{pl}_{\,i_v@x\leftrightarrow i_v@x'}
 \bigl(E_h(x),E_h(x'),
       h[x^0][x'^{\,0}]\cdot\operatorname{Tag}_0(e)\bigr).
\]
The final premise has type
$v\diamond_m(x^0)\leftrightarrow v\diamond_m(x'^{\,0})$,
so this is a proof of $\operatorname{Ext}(i_v)$.

Fix $x:X$ and write
\[
 P_x=\operatorname{Pair}(\operatorname{Part}(v),x^0),
 \qquad
 i_{P_x}=\operatorname{pred}_X(x'.\,P_x@(x'^{\,\ell})).
\]
The Pair laws, under the present hypothesis $h$, supply
\begin{align*}
 L_h(x')&:P_x@(x'^{\,\ell})
                  \leftrightarrow v\diamond_m(x'^{\,0}),\\
 R_h(x')&:P_x@(x'^{\,r})\leftrightarrow x\simeq x'.
\end{align*}
At an arbitrary $x':X$, denote the two predicate evaluation proofs by
\begin{align*}
 r_{P_x}(x')&:i_{P_x}@x'
              \leftrightarrow P_x@((\operatorname{Round}_Xx')^\ell),\\
 r_v(x')&:i_v@x'
              \leftrightarrow v\diamond_m((\operatorname{Round}_Xx')^0).
\end{align*}
Using $L_h(\operatorname{Round}_Xx')$ between them gives
\[
 d=\Lambda x'.\,
      \mathsf{pl}_{\,i_{P_x}@x'\leftrightarrow i_v@x'}
      (r_{P_x}(x'),r_v(x'),
       L_h(\operatorname{Round}_Xx'))
      :i_{P_x}\approx i_v.
\]
Consequently
\[
 ex_P=\operatorname{back}(\operatorname{ExtC}(d),ex_v)
       :\operatorname{Ext}(i_{P_x}).
\]

We first construct both directions of
\begin{equation}
 \mathcal P_R(P_x)\leftrightarrow\operatorname{le}(i_v,x),
 \label{rel-r-property-law}
\end{equation}
where, by the definition of $\mathcal P_R$,
\[
 \mathcal P_R(P_x)=\operatorname{Ext}(i_{P_x})\Rightarrow
 \forall x':X.\bigl(\operatorname{Valid}(x')\Rightarrow
 (P_x@(x'^{\,r})\Rightarrow\operatorname{le}(i_{P_x},x'))\bigr).
\]
For the forward direction, assume $g:\mathcal P_R(P_x)$ and put
\[
 p_x=\operatorname{back}(R_h(x),\operatorname{refl}(x))
           :P_x@(x^r).
\]
Then
\[
 (g\cdot ex_P)[x]\cdot u\cdot p_x
      :\operatorname{le}(i_{P_x},x),
 \qquad
 u[i_{P_x}][i_v]\cdot d:
      \operatorname{le}(i_{P_x},x)\leftrightarrow
      \operatorname{le}(i_v,x).
\]
Thus the forward builder is
\[
 \lambda^*g.\,
 \operatorname{to}\bigl(u[i_{P_x}][i_v]\cdot d,
                      (g\cdot ex_P)[x]\cdot u\cdot p_x\bigr).
\]
For the reverse direction, assume $l:\operatorname{le}(i_v,x)$.
Introduce $ex':\operatorname{Ext}(i_{P_x})$, $x':X$,
$u':\operatorname{Valid}(x')$, and $p:P_x@(x'^{\,r})$.
The proof $ex'$ need not be used.  Set
\[
 e=\operatorname{to}(R_h(x'),p):x\simeq x'.
\]
Here $e[i_v]:\operatorname{le}(i_v,x)\leftrightarrow
\operatorname{le}(i_v,x')$, whereas
\[
 u'[i_{P_x}][i_v]\cdot d:
 \operatorname{le}(i_{P_x},x')\leftrightarrow
 \operatorname{le}(i_v,x').
\]
The required proof of $\operatorname{le}(i_{P_x},x')$ is therefore
\[
 \operatorname{back}\bigl(u'[i_{P_x}][i_v]\cdot d,
                         \operatorname{to}(e[i_v],l)\bigr).
\]
Abstracting, in order, over $p,u',x',ex',l$ gives the reverse
direction of \eqref{rel-r-property-law}; pairing the directions gives
the displayed biconditional.

Finally, by definition,
$\operatorname{sb}(v,m,x^0)=\operatorname{qry}_R(m,P_x)$.
The query comparison is
\[
 q_R(P_x):\operatorname{qry}_R(m,P_x)
                         \leftrightarrow\mathcal P_R(P_x).
\]
Composing its forward direction with the forward direction of
\eqref{rel-r-property-law}, and composing the two reverse directions
in the reverse order, gives \eqref{rel-sb-law}.
\end{proof}

\begin{lemma}[The shift and its predecessor comparison]
\label{rel-shift}
Define the data operations
\[
 \delta x=\operatorname{lift}_m(x^0):X,
 \qquad
 Si=\operatorname{pred}_X(x.\,i@\delta x):\mathsf P_X.
\]
For every $x:X$ there is a proof
$V_\delta(x):\operatorname{Valid}(\delta x)$, and for
$e:x\simeq x'$ there is a proof
$\operatorname{DeltaC}(e):\delta x\simeq\delta x'$.
For $ex:\operatorname{Ext}(i)$ there are builders
\begin{align}
 s_{ex}(x)&:Si@x\leftrightarrow i@\delta x,
                                    \label{rel-shift-evaluation}\\
 \operatorname{ExtS}(i,ex)&:\operatorname{Ext}(Si),
                                    \label{rel-shift-extensionality}\\
 d_s&:i_{z_i}\approx Si.             \label{rel-shift-comparison}
\end{align}
For the additional input $u:\operatorname{Valid}(x)$ there is a builder
\begin{equation}
 a_{ex,u}:\operatorname{le}(i,\delta x)
                    \leftrightarrow\operatorname{le}(Si,x).
 \label{rel-shift-predecessor}
\end{equation}
\end{lemma}

\begin{proof}
The established validity proof for a lift, applied to $x^0$, is
$V_\delta(x)$.  Congruence is the composite
\[
 \operatorname{DeltaC}(e)
       =\operatorname{LC}(\operatorname{Tag}_0(e)).
\]
Neither construction requires $x$ to be valid or any predicate to
be extensional.

The predicate evaluation law gives
$Si@x\leftrightarrow i@\delta(\operatorname{Round}_Xx)$.
The proof
\[
 ex[\delta x][\delta(\operatorname{Round}_Xx)]\cdot
      \operatorname{DeltaC}(\operatorname{RoundRel}_X(x))
 :i@\delta x\leftrightarrow
      i@\delta(\operatorname{Round}_Xx)
\]
and that evaluation give $s_{ex}(x)$ by $\mathsf{pl}$ at the target
in \eqref{rel-shift-evaluation}.  Set
\[
 \operatorname{ExtS}(i,ex)=
 \Lambda x,x'.\lambda^*e.\,
 \mathsf{pl}_{\,Si@x\leftrightarrow Si@x'}
 \bigl(s_{ex}(x),s_{ex}(x'),
       ex[\delta x][\delta x']\cdot\operatorname{DeltaC}(e)\bigr).
\]
This has the type in \eqref{rel-shift-extensionality}.

From $ex$, the earlier response builder gives
$h_i:\operatorname{Resp}(z_i)$, and the earlier exact evaluation
gives, for every $y:L$, a proof
\[
 Z_{ex}(y):z_i\diamond_m y\leftrightarrow
                i@\operatorname{lift}_m(y).
\]
In particular, at $y=x^0$ the right-hand side is $i@\delta x$.
Together with $E_{h_i}(x)$ and $s_{ex}(x)$ this gives
\[
 d_s=\Lambda x.\,
 \mathsf{pl}_{\,i_{z_i}@x\leftrightarrow Si@x}
 \bigl(E_{h_i}(x),Z_{ex}(x^0),s_{ex}(x)\bigr).
\]

Now assume $u:\operatorname{Valid}(x)$.  The three available
biconditionals, with their exact endpoints, are
\begin{align*}
 \operatorname{LP}(i,x^0)&:
 \operatorname{le}(i,\delta x)
       \leftrightarrow\operatorname{sb}(z_i,m,x^0),\\
 \operatorname{SB}(z_i,h_i,x,u)&:
 \operatorname{sb}(z_i,m,x^0)
       \leftrightarrow\operatorname{le}(i_{z_i},x),\\
 u[i_{z_i}][Si]\cdot d_s&:
 \operatorname{le}(i_{z_i},x)
       \leftrightarrow\operatorname{le}(Si,x).
\end{align*}
The forward direction of $a_{ex,u}$ is
\[
 \lambda^*t.\,
 \operatorname{to}\bigl(u[i_{z_i}][Si]\cdot d_s,
  \operatorname{to}\bigl(\operatorname{SB}(z_i,h_i,x,u),
   \operatorname{to}(\operatorname{LP}(i,x^0),t)\bigr)\bigr).
\]
Its reverse direction is
\[
 \lambda^*t.\,
 \operatorname{back}\bigl(\operatorname{LP}(i,x^0),
  \operatorname{back}\bigl(\operatorname{SB}(z_i,h_i,x,u),
   \operatorname{back}(u[i_{z_i}][Si]\cdot d_s,t)\bigr)\bigr).
\]
Pairing them proves \eqref{rel-shift-predecessor}.  The only use of
$u$ is in the response comparison and in transport at the fixed
argument $x$; no unconditional replacement of $i_{z_i}$ by $Si$
inside $\operatorname{le}$ has been made.
\end{proof}

\begin{lemma}[Induction respects predicate equivalence]
\label{rel-ind-congruence}
For $d:i\approx i'$ there is a builder
\[
 \operatorname{IndC}(d):
       \operatorname{Ind}(i)\leftrightarrow\operatorname{Ind}(i').
\]
This builder needs no extensionality hypothesis on either predicate.
\end{lemma}

\begin{proof}
For $x:X$ and $u:\operatorname{Valid}(x)$ put
\[
 c_u=u[i][i']\cdot d:
       \operatorname{le}(i,x)\leftrightarrow\operatorname{le}(i',x).
\]
The two directions, paired in the displayed order, are
\begin{align*}
 &\lambda^*l.\Lambda x.\lambda^*u.\lambda^*t'.\,
    \operatorname{to}\bigl(d[x],
       l[x]\cdot u\cdot\operatorname{back}(c_u,t')\bigr),\\
 &\lambda^*l'.\Lambda x.\lambda^*u.\lambda^*t.\,
    \operatorname{back}\bigl(d[x],
       l'[x]\cdot u\cdot\operatorname{to}(c_u,t)\bigr).
\end{align*}
In the first line $l:\operatorname{Ind}(i)$ and
$t':\operatorname{le}(i',x)$; in the second
$l':\operatorname{Ind}(i')$ and $t:\operatorname{le}(i,x)$.
Thus every application has the exact antecedent required by
\eqref{rel-ind-definition}.
\end{proof}

\begin{lemma}[A valid induction object]
\label{rel-wf}
Define, without proof hypotheses,
\[
 \operatorname{WF}=\operatorname{mk}_X(v.\,\operatorname{Ind}(i_v)):X.
\]
There is a proof $V_{\operatorname{WF}}:
\operatorname{Valid}(\operatorname{WF})$.  For
$ex:\operatorname{Ext}(i)$ there is a builder
\begin{equation}
 w_{ex}:\operatorname{le}(i,\operatorname{WF})
                       \leftrightarrow\operatorname{Ind}(Si).
 \label{rel-wf-evaluation}
\end{equation}
Finally, from $i:\mathsf P_X$, $ex:\operatorname{Ext}(i)$ and
$l:\operatorname{Ind}(i)$ there is a builder
\begin{equation}
 b(i,ex,l):i@\operatorname{WF}.
 \label{rel-induction-builder}
\end{equation}
\end{lemma}

\begin{proof}
Since $\operatorname{le}(i,\operatorname{WF})=[\operatorname{WF}z_i]$,
the data cancellation law gives, for every $i$, a proof
\[
 F(i):\operatorname{le}(i,\operatorname{WF})
                   \leftrightarrow\operatorname{Ind}(i_{z_i}).
\]
This evaluation uses neither $\operatorname{Resp}(z_i)$ nor
$\operatorname{Ext}(i)$.

To prove validity, take arbitrary $i,i':\mathsf P_X$ and
$d:i\approx i'$.  We construct
\[
 D(d):i_{z_i}\approx i_{z_{i'}}
\]
without introducing extensionality assumptions.  At $x:X$ put
$t=(\operatorname{Round}_Xx)^0:L$.  The predicate evaluation law
and the earlier rounded evaluation of $z_i$ give the following
chain of biconditionals:
\[
 i_{z_i}@x
 \ \leftrightarrow\ z_i\diamond_m t
 \ \leftrightarrow\ i@\operatorname{lift}_m(\operatorname{Round}t)
 \ \leftrightarrow\ i'@\operatorname{lift}_m(\operatorname{Round}t)
 \ \leftrightarrow\ z_{i'}\diamond_m t
 \ \leftrightarrow\ i_{z_{i'}}@x.
\]
The middle comparison is exactly
$d[\operatorname{lift}_m(\operatorname{Round}t)]$.
The other four comparisons are the two predicate evaluations and
the two rounded evaluations.  Applying $\mathsf{pl}$ to those five
proofs at target $i_{z_i}@x\leftrightarrow i_{z_{i'}}@x$, then
abstracting over $x$, constructs $D(d)$.  Therefore
\[
 V_{\operatorname{WF}}=
 \Lambda i,i'.\lambda^*d.\,
 \mathsf{pl}_{\,\operatorname{le}(i,\operatorname{WF})
                   \leftrightarrow
                   \operatorname{le}(i',\operatorname{WF})}
 \bigl(F(i),F(i'),\operatorname{IndC}(D(d))\bigr).
\]

Now assume $ex:\operatorname{Ext}(i)$.  The proof $d_s$ from
\eqref{rel-shift-comparison} and Lemma~\ref{rel-ind-congruence} give
\[
 \operatorname{IndC}(d_s):
          \operatorname{Ind}(i_{z_i})\leftrightarrow\operatorname{Ind}(Si).
\]
The forward and reverse directions of $w_{ex}$ are, respectively,
\begin{align*}
 &\lambda^*t.\,
   \operatorname{to}(\operatorname{IndC}(d_s),
                     \operatorname{to}(F(i),t)),\\
 &\lambda^*l_S.\,
   \operatorname{back}(F(i),
                       \operatorname{back}(\operatorname{IndC}(d_s),l_S)).
\end{align*}
Their pair proves \eqref{rel-wf-evaluation}.

For the final assertion, fix $l:\operatorname{Ind}(i)$.  We first
construct $l_S:\operatorname{Ind}(Si)$.  Under $x:X$,
$u:\operatorname{Valid}(x)$ and $t:\operatorname{le}(Si,x)$,
the proof $\operatorname{back}(a_{ex,u},t)$ has type
$\operatorname{le}(i,\delta x)$.  Since
$V_\delta(x):\operatorname{Valid}(\delta x)$, it follows that
\[
 l[\delta x]\cdot V_\delta(x)\cdot
                   \operatorname{back}(a_{ex,u},t):i@\delta x.
\]
Applying the reverse direction of $s_{ex}(x)$ gives $Si@x$.
Thus the required builder is
\[
 l_S=\Lambda x.\lambda^*u.\lambda^*t.\,
 \operatorname{back}\bigl(s_{ex}(x),
       l[\delta x]\cdot V_\delta(x)\cdot
                   \operatorname{back}(a_{ex,u},t)\bigr).
\]
Since $\operatorname{back}(w_{ex},l_S)$ has type
$\operatorname{le}(i,\operatorname{WF})$, define
\[
 b(i,ex,l)=
 l[\operatorname{WF}]\cdot V_{\operatorname{WF}}\cdot
                      \operatorname{back}(w_{ex},l_S).
\]
Its type is $i@\operatorname{WF}$, as required.
\end{proof}

\begin{proposition}[The diagonal contradiction]
\label{rel-contradiction}
The specialized relational construction yields a proof of
$\bot=F_b$ in the ambient labelled context.
\end{proposition}

\begin{proof}
Define the data formula and data predicate
\begin{align*}
 Q(x)&=\forall i:\mathsf P_X.\bigl(\operatorname{Ext}(i)
           \Rightarrow(\operatorname{le}(i,x)\Rightarrow i@\delta x)\bigr),\\
 j_*&=\operatorname{pred}_X(x.\,\neg Q(x)):\mathsf P_X.
\end{align*}
Neither definition uses proof variables.

First we construct
\[
 \operatorname{QC}(e):Q(x)\leftrightarrow Q(x')
 \qquad(e:x\simeq x').
\]
For any $i:\mathsf P_X$ and $ex:\operatorname{Ext}(i)$, put
\[
 c_{ex,e}=ex[\delta x][\delta x']\cdot\operatorname{DeltaC}(e)
               :i@\delta x\leftrightarrow i@\delta x'.
\]
The forward direction of $\operatorname{QC}(e)$ is
\[
 \lambda^*q.\Lambda i.\lambda^*ex.\lambda^*t'.\,
 \operatorname{to}\bigl(c_{ex,e},
       q[i]\cdot ex\cdot\operatorname{back}(e[i],t')\bigr),
\]
where $q:Q(x)$ and $t':\operatorname{le}(i,x')$.
Its reverse direction is
\[
 \lambda^*q'.\Lambda i.\lambda^*ex.\lambda^*t.\,
 \operatorname{back}\bigl(c_{ex,e},
       q'[i]\cdot ex\cdot\operatorname{to}(e[i],t)\bigr),
\]
where $q':Q(x')$ and $t:\operatorname{le}(i,x)$.
Pairing these terms gives the asserted congruence.

The predicate evaluation law supplies
$j_*@x\leftrightarrow\neg Q(\operatorname{Round}_Xx)$.
Using $\operatorname{QC}(\operatorname{RoundRel}_X(x))$ and
$\mathsf{pl}$ with that evaluation yields
\begin{equation}
 n(x):j_*@x\leftrightarrow\neg Q(x).
 \label{rel-negative-evaluation}
\end{equation}
The extensionality builder is
\[
 ex_*=
 \Lambda x,x'.\lambda^*e.\,
 \mathsf{pl}_{\,j_*@x\leftrightarrow j_*@x'}
       \bigl(n(x),n(x'),\operatorname{QC}(e)\bigr)
       :\operatorname{Ext}(j_*).
\]
In both uses of $\mathsf{pl}$ the passage from a biconditional
between the $Q$-formulas to the corresponding negated formulas is
propositional; no predicate is substituted inside data.

We next construct $l_*:\operatorname{Ind}(j_*)$.
Introduce $x:X$, $u:\operatorname{Valid}(x)$ and
$t:\operatorname{le}(j_*,x)$.  To obtain $j_*@x$, it suffices,
by the reverse direction of $n(x)$, to obtain $\neg Q(x)$.
Under the further assumption $q:Q(x)$ we have
\[
 q[j_*]\cdot ex_*\cdot t:j_*@\delta x,
\]
and hence
\begin{equation}
 n_\delta=
 \operatorname{to}\bigl(n(\delta x),q[j_*]\cdot ex_*\cdot t\bigr)
                  :\neg Q(\delta x).
 \label{rel-negative-shift}
\end{equation}
To construct the opposite formula $Q(\delta x)$, take arbitrary
$i:\mathsf P_X$, $ex:\operatorname{Ext}(i)$ and
$t':\operatorname{le}(i,\delta x)$.  At this point
$a_{ex,u}$ uses the original proof $u:\operatorname{Valid}(x)$.
Its forward direction gives
\[
 \operatorname{to}(a_{ex,u},t'):\operatorname{le}(Si,x).
\]
Because $\operatorname{ExtS}(i,ex):\operatorname{Ext}(Si)$,
we may instantiate $q$ at $Si$ to obtain
\[
 q[Si]\cdot\operatorname{ExtS}(i,ex)\cdot
           \operatorname{to}(a_{ex,u},t'):Si@\delta x.
\]
The forward direction of $s_{ex}(\delta x)$ turns this into
$i@\delta(\delta x)$.  Therefore
\begin{equation}
 \begin{split}
 p_\delta={}&\Lambda i.\lambda^*ex.\lambda^*t'.\,
 \operatorname{to}\bigl(s_{ex}(\delta x),\\[-2pt]
 &\hspace{35mm}
 q[Si]\cdot\operatorname{ExtS}(i,ex)\cdot
                   \operatorname{to}(a_{ex,u},t')\bigr)
       :Q(\delta x).
 \end{split}
 \label{rel-positive-shift}
\end{equation}
Combining \eqref{rel-negative-shift} and \eqref{rel-positive-shift}
gives $n_\delta\cdot p_\delta:\bot$.  With these local terms
inlined in their indicated context, define
\[
 l_*=
 \Lambda x.\lambda^*u.\lambda^*t.\,
 \operatorname{back}\bigl(n(x),
                  \lambda^*q.\,n_\delta\cdot p_\delta\bigr).
\]
Its body has type $j_*@x$, so
$l_*:\operatorname{Ind}(j_*)$.

Finally construct $q_0:Q(\operatorname{WF})$.  Under arbitrary
$i:\mathsf P_X$, $ex:\operatorname{Ext}(i)$ and
$t:\operatorname{le}(i,\operatorname{WF})$, the forward direction
of $w_{ex}$ gives $\operatorname{Ind}(Si)$.  Applying the induction
builder to the predicate $Si$ therefore gives
\[
 b\bigl(Si,\operatorname{ExtS}(i,ex),
                 \operatorname{to}(w_{ex},t)\bigr):Si@\operatorname{WF}.
\]
Here the builder $b$ is instantiated at $Si$, while the displayed
$w_{ex}$ and the following $s_{ex}$ are the builders for the
original predicate $i$.  Thus
\[
 q_0=\Lambda i.\lambda^*ex.\lambda^*t.\,
 \operatorname{to}\bigl(s_{ex}(\operatorname{WF}),
       b(Si,\operatorname{ExtS}(i,ex),
                        \operatorname{to}(w_{ex},t))\bigr)
       :Q(\operatorname{WF}).
\]
On the other hand,
$b(j_*,ex_*,l_*):j_*@\operatorname{WF}$, and hence
\[
 \operatorname{to}\bigl(n(\operatorname{WF}),
                      b(j_*,ex_*,l_*)\bigr)
                  :\neg Q(\operatorname{WF}).
\]
The required ambient proof is therefore
\[
 \operatorname{to}\bigl(n(\operatorname{WF}),
                      b(j_*,ex_*,l_*)\bigr)\cdot q_0:\bot.
\]
\end{proof}

\begin{remark}[Dependencies and discharged assumptions]
\label{rel-diagonal-dependencies}
At the fixed specialization, $i_v$ is defined before its
response and extensionality proofs; $\delta$ and $S$ are data
definitions independent of those proofs; $\operatorname{WF}$ uses
the already defined $i_v$ and $\operatorname{Ind}$; and $Q,j_*$ use
only the previously defined data operations.  In particular,
$\operatorname{WF}$ is defined even when $v$ is not responsive.
Its validity proof uses rounded evaluations and does not assume
extensionality of the predicates it quantifies over.

The inputs $h,ex,u$ occur only in proof builders.  Every use of
\eqref{rel-sb-law} or \eqref{rel-shift-predecessor} supplies the
stated response, extensionality and validity inputs.  Every logical
quantifier introduced by the displayed tail builders ranges over
$X$ or $\mathsf P_X$ with its fixed direct-domain annotation.
After inlining the local builders,
the proof of Proposition~\ref{rel-contradiction} has no free proof
or data assumptions beyond the ambient context: $A_0,m_0,ad$ are
the earlier constructed specialization, and every introduced
logical assumption has been abstracted or supplied as an argument.
\end{remark}
